\documentclass[11pt]{article}
\usepackage[margin=0.8in]{geometry}
\usepackage[T1]{fontenc}
\usepackage{lmodern}
\usepackage{textcomp}
\usepackage{listings}

\lstset{
  basicstyle=\ttfamily\small,
  columns=fullflexible,
  keepspaces=true,
  showstringspaces=false,
  upquote=true,
  breaklines=true,
  aboveskip=3pt,
  belowskip=3pt
}
\lstdefinestyle{answer}{frame=single, framesep=5pt}
\renewcommand{\labelenumi}{(\alph{enumi})}
\setlength{\parindent}{0pt}

% EDITING: Replace TODO inside each answer listing with ordinary Haskell.
% Backslashes, arrows, underscores, and multi-line answers need no escaping
% inside lstlisting. For example, write a lambda as: \x -> ...
% Compile this file again after editing to update its PDF.
% Original practice questions modeled on Q2 and Q3 of the 2024 midterm.
% No answer key is included.

\begin{document}
\raggedbottom
\section*{Haskell Types Practice: October 8, 2026}
This worksheet follows Questions 2 and 3 of the Spring 2024 CMSC 433
midterm, with different expressions and target types. The two main
sections are worth 20 points each. Optional challenges follow.

\textbf{Editing:} Replace each \texttt{TODO} in this source file with
your answer. Answer boxes accept ordinary Haskell syntax.

\subsection*{Question 2: What is the type? (20 points)}
For each expression, give its \textbf{most general Haskell type}, including
all necessary typeclass constraints, or write \texttt{ill-typed} and
briefly explain the mismatch. \textbf{Exam convention: treat all raw
numeric literals as \texttt{Int}.} Give the most general type under
that convention. Use the reference appendix for library signatures.

\begin{enumerate}
\item
\begin{minipage}[t]{\linewidth}
\vspace{-\baselineskip}
\begin{lstlisting}
\x y -> (y, x, x)
\end{lstlisting}
\begin{lstlisting}[style=answer]
a -> b -> (b, a, a)
\end{lstlisting}
\end{minipage}

\item
\begin{minipage}[t]{\linewidth}
\vspace{-\baselineskip}
\begin{lstlisting}
\x y -> if x < y then show x else show (y, x)
\end{lstlisting}
\begin{lstlisting}[style=answer]
(Ord a, Show a) => a -> a -> String
\end{lstlisting}
\end{minipage}

\item
\begin{minipage}[t]{\linewidth}
\vspace{-\baselineskip}
\begin{lstlisting}
\x xs -> x : reverse xs ++ [x]
\end{lstlisting}
\begin{lstlisting}[style=answer]
a -> [a] -> [a]
\end{lstlisting}
\end{minipage}

\item
\begin{minipage}[t]{\linewidth}
\vspace{-\baselineskip}
\begin{lstlisting}
foldr (const False)
\end{lstlisting}
\begin{lstlisting}[style=answer]
ill typed
\end{lstlisting}
\end{minipage}

\item
\begin{minipage}[t]{\linewidth}
\vspace{-\baselineskip}
\begin{lstlisting}
(, 'q')
\end{lstlisting}
\begin{lstlisting}[style=answer]
a -> (a, Char)
\end{lstlisting}
\end{minipage}

\item
\begin{minipage}[t]{\linewidth}
\vspace{-\baselineskip}
\begin{lstlisting}
(,) ['q']
\end{lstlisting}
\begin{lstlisting}[style=answer]
b -> ([Char], b)
\end{lstlisting}
\end{minipage}

\item
\begin{minipage}[t]{\linewidth}
\vspace{-\baselineskip}
\begin{lstlisting}
reverse . map Just
\end{lstlisting}
\begin{lstlisting}[style=answer]
[a] -> [Maybe a]
\end{lstlisting}
\end{minipage}

\item
\begin{minipage}[t]{\linewidth}
\vspace{-\baselineskip}
\begin{lstlisting}
\xs -> filter (>= 7) (reverse xs)
\end{lstlisting}
\begin{lstlisting}[style=answer]
[Int] -> [Int]
\end{lstlisting}
\end{minipage}

\item
\begin{minipage}[t]{\linewidth}
\vspace{-\baselineskip}
\begin{lstlisting}
let keep x = x in (keep 'q', keep [False])
\end{lstlisting}
\begin{lstlisting}[style=answer]
(Char, [Bool])
\end{lstlisting}
\end{minipage}

\item
\begin{minipage}[t]{\linewidth}
\vspace{-\baselineskip}
\begin{lstlisting}
fmap (fmap (== 'q')) [Nothing, Just 'r']
\end{lstlisting}
\begin{lstlisting}[style=answer]
[Maybe Bool]
\end{lstlisting}
\end{minipage}

\end{enumerate}

\subsection*{Question 3: Give an expression matching the type (20 points)}
For each type below, write a Haskell expression that has that type.
Do not write trivial expressions (such as \texttt{[]}, \texttt{Nothing},
or \texttt{undefined}) unless there is no other option. You may use any
function from the appendix, \texttt{do} syntax, list comprehensions,
or any valid Haskell.

\begin{enumerate}
\item
\begin{minipage}[t]{\linewidth}
\vspace{-\baselineskip}
\begin{lstlisting}
Bool -> Maybe Bool
\end{lstlisting}
\begin{lstlisting}[style=answer]
\x -> if x then Just x else Nothing


\end{lstlisting}
\end{minipage}

\item
\begin{minipage}[t]{\linewidth}
\vspace{-\baselineskip}
\begin{lstlisting}
a -> [b]
\end{lstlisting}
\begin{lstlisting}[style=answer]
\x -> []


\end{lstlisting}
\end{minipage}

\item
\begin{minipage}[t]{\linewidth}
\vspace{-\baselineskip}
\begin{lstlisting}
(Char -> Int -> Bool) -> [Char] -> [Int] -> [Bool]
\end{lstlisting}
\begin{lstlisting}[style=answer]
liftA2


\end{lstlisting}
\end{minipage}

\item
\begin{minipage}[t]{\linewidth}
\vspace{-\baselineskip}
\begin{lstlisting}
(a -> b -> c) -> Maybe a -> Maybe b -> Maybe (c, b)
\end{lstlisting}
\begin{lstlisting}[style=answer]
\f -> liftA2 (\x y -> (f x y, y))


\end{lstlisting}
\end{minipage}

\item
\begin{minipage}[t]{\linewidth}
\vspace{-\baselineskip}
\begin{lstlisting}
(a -> b) -> (a -> Bool) -> [a] -> [b]
\end{lstlisting}
\begin{lstlisting}[style=answer]
\f f2 xs -> map f (filter f2 xs)


\end{lstlisting}
\end{minipage}

\item
\begin{minipage}[t]{\linewidth}
\vspace{-\baselineskip}
\begin{lstlisting}
Maybe a -> (a -> [b]) -> [(b, a)]
\end{lstlisting}
\begin{lstlisting}[style=answer]
\ma f -> maybe [] (\x -> map (\y -> (y, x)) (f x)) ma


\end{lstlisting}
\end{minipage}

\item
\begin{minipage}[t]{\linewidth}
\vspace{-\baselineskip}
\begin{lstlisting}
Eq a => [a] -> a -> [Bool]
\end{lstlisting}
\begin{lstlisting}[style=answer]
\(x:xs) y -> if y == x then [True] else [False]


\end{lstlisting}
\end{minipage}

\item
\begin{minipage}[t]{\linewidth}
\vspace{-\baselineskip}
\begin{lstlisting}
Show a => Maybe a -> IO String
\end{lstlisting}
\begin{lstlisting}[style=answer]
\ma -> return (show ma)



\end{lstlisting}
\end{minipage}

\item
\begin{minipage}[t]{\linewidth}
\vspace{-\baselineskip}
\begin{lstlisting}
(a -> b -> c) -> (b, a) -> c
\end{lstlisting}
\begin{lstlisting}[style=answer]
\f (b, a) -> f a b 


\end{lstlisting}
\end{minipage}

\item
\begin{minipage}[t]{\linewidth}
\vspace{-\baselineskip}
\begin{lstlisting}
((a, b) -> c) -> b -> c
\end{lstlisting}
\begin{lstlisting}[style=answer]
undefined


\end{lstlisting}
\end{minipage}

\end{enumerate}

\subsection*{Optional challenges}
These are additional practice and are not part of the 40-point worksheet.

\begin{enumerate}
\item
\begin{minipage}[t]{\linewidth}
Give the most general type or explain the type error. Compare
the reason for your answer with Question 2(i).
\begin{lstlisting}
\keep -> (keep 'q', keep [False])
\end{lstlisting}
\begin{lstlisting}[style=answer]
TODO

\end{lstlisting}
\end{minipage}

\item
\begin{minipage}[t]{\linewidth}
Give the most general type or explain the type error. Identify
the structures acted on by both occurrences of \texttt{fmap}.
\begin{lstlisting}
fmap (fmap not) [Just True, Nothing, Just 'q']
\end{lstlisting}
\begin{lstlisting}[style=answer]
TODO

\end{lstlisting}
\end{minipage}

\item
\begin{minipage}[t]{\linewidth}
Give the most general type. Which type variable belongs to the
input container, and which belongs to the accumulator?
\begin{lstlisting}
foldl (flip (:)) []
\end{lstlisting}
\begin{lstlisting}[style=answer]
TODO

\end{lstlisting}
\end{minipage}

\item
\begin{minipage}[t]{\linewidth}
Give an expression of this type. Apply the function to each
input, discard absent results, and pair each present result with its
original input. Preserve input order, including duplicates.
\begin{lstlisting}
(a -> Maybe b) -> [a] -> [(a, b)]
\end{lstlisting}
\begin{lstlisting}[style=answer]
TODO



\end{lstlisting}
\end{minipage}

\end{enumerate}

\subsection*{Reference appendix}
These signatures use the generalized \texttt{Foldable} versions from
the old exam's appendix. \texttt{String} is a synonym for \texttt{[Char]}.
Tuple sections, such as \texttt{(, False)}, are allowed as in the exam;
enable \texttt{TupleSections} if checking them in GHCi.

\begin{lstlisting}
data Bool = False | True
data Maybe a = Nothing | Just a

id :: a -> a
const :: a -> b -> a
flip :: (a -> b -> c) -> b -> a -> c
(.) :: (b -> c) -> (a -> b) -> a -> c
(,) :: a -> b -> (a, b)

not :: Bool -> Bool
(==), (/=) :: Eq a => a -> a -> Bool
(<), (<=), (>), (>=) :: Ord a => a -> a -> Bool
(+), (-), (*) :: Num a => a -> a -> a
show :: Show a => a -> String

(:) :: a -> [a] -> [a]
(++) :: [a] -> [a] -> [a]
reverse :: [a] -> [a]
map :: (a -> b) -> [a] -> [b]
filter :: (a -> Bool) -> [a] -> [a]
foldr :: Foldable t => (a -> b -> b) -> b -> t a -> b
foldl :: Foldable t => (b -> a -> b) -> b -> t a -> b

fmap :: Functor f => (a -> b) -> f a -> f b
pure :: Applicative f => a -> f a
(<*>) :: Applicative f => f (a -> b) -> f a -> f b
liftA2 :: Applicative f => (a -> b -> c) -> f a -> f b -> f c
return :: Monad m => a -> m a
(>>=) :: Monad m => m a -> (a -> m b) -> m b
(>>) :: Monad m => m a -> m b -> m b

maybe :: b -> (a -> b) -> Maybe a -> b
maybeToList :: Maybe a -> [a]
mapMaybe :: (a -> Maybe b) -> [a] -> [b]

print :: Show a => a -> IO ()
putStrLn :: String -> IO ()
\end{lstlisting}
\textbf{Numeric literals:} For this worksheet, raw numbers have type
\texttt{Int}, following the TA's exam clarification. When a type is
fixed to \texttt{Int}, omit constraints such as \texttt{Ord Int} and
\texttt{Num Int}, since those instances are already known. Library
signatures above remain general. In GHCi, annotate a literal with
\texttt{:: Int} to reproduce the exam convention when checking a type.

\textbf{Constraints:} \texttt{Ord a} implies \texttt{Eq a}; it does not
imply \texttt{Num a} or \texttt{Show a}. Lists and \texttt{Maybe} are
instances of \texttt{Functor}, \texttt{Applicative}, and \texttt{Monad}.
They are also \texttt{Foldable}. Tuples and \texttt{Maybe a} have
\texttt{Show} instances when their contained types have \texttt{Show}
instances.

\end{document}
